% =============================================================================
%
%  TDK Thesis - University of Debrecen
%  Graph-Structured Adaptive Memory: Evolving Contexts for Agentic,
%  Self-Improving Systems
%
%  Author: Machungo Erick
%  Date: 2026
%
% =============================================================================

\documentclass[12pt,a4paper,oneside]{article}

% =============================================================================
% PAGE LAYOUT (TDK Requirements)
% =============================================================================
\usepackage[left=3cm, right=2cm, top=3cm, bottom=3cm]{geometry}

% =============================================================================
% FONTS AND ENCODING
% =============================================================================
\usepackage{mathptmx}           % Times New Roman equivalent
\usepackage{xspace}             % Smart spacing for macros
\usepackage[T1]{fontenc}        % Proper font encoding for accented characters
\usepackage[utf8]{inputenc}
\usepackage{textcomp}           % Additional text symbols
\usepackage{newunicodechar}     % Support for custom Unicode characters
\newunicodechar{ő}{\H{o}}       % Hungarian ő
\newunicodechar{ű}{\H{u}}       % Hungarian ű

% =============================================================================
% SPACING
% =============================================================================
\usepackage{setspace}
\onehalfspacing                 % 1.5 line spacing per TDK requirements

% =============================================================================
% COLORS (define before hyperref)
% =============================================================================
\usepackage{xcolor}

% Custom color palette
\definecolor{thesisblue}{RGB}{0, 83, 156}       % University blue
\definecolor{linkblue}{RGB}{0, 102, 204}        % Hyperlink blue
\definecolor{citegreen}{RGB}{0, 128, 0}         % Citation green
\definecolor{codegray}{RGB}{245, 245, 245}      % Code background
\definecolor{codetext}{RGB}{40, 40, 40}         % Code text

% =============================================================================
% GRAPHICS AND FIGURES
% =============================================================================
\usepackage{graphicx}
\graphicspath{{}}
\usepackage{subcaption}
\usepackage{float}

% =============================================================================
% TABLES
% =============================================================================
\usepackage{booktabs}
\usepackage{multirow}
\usepackage{array}
\usepackage{tabularx}

% =============================================================================
% MATHEMATICS
% =============================================================================
\usepackage{amsmath,amssymb,amsfonts}
\usepackage{amsthm}

\newtheorem{proposition}{Proposition}
\newtheorem{definition}{Definition}

% =============================================================================
% ALGORITHMS
% =============================================================================
\usepackage{algorithm}
\usepackage{algorithmic}

% =============================================================================
% CITATIONS
% =============================================================================
\usepackage[authoryear,round]{natbib}
\bibliographystyle{plainnat}

% =============================================================================
% CODE LISTINGS
% =============================================================================
\usepackage{listings}
\lstset{
    basicstyle=\small\ttfamily\color{codetext},
    breaklines=true,
    frame=single,
    framerule=0.5pt,
    rulecolor=\color{gray!50},
    backgroundcolor=\color{codegray},
    captionpos=b,
    numbers=left,
    numberstyle=\tiny\color{gray},
    numbersep=8pt,
    tabsize=4,
    showstringspaces=false
}

% =============================================================================
% LISTS
% =============================================================================
\usepackage{enumitem}
\setlist{nosep, leftmargin=*}
\setlist[itemize]{label=\textbullet}
\setlist[enumerate]{label=\arabic*.}

% =============================================================================
% SECTION FORMATTING (14pt sections, 13pt subsections per TDK)
% =============================================================================
\usepackage{titlesec}

\titleformat{\section}
    {\fontsize{14}{17}\bfseries}
    {\thesection}
    {1em}
    {}

\titleformat{\subsection}
    {\fontsize{13}{16}\bfseries}
    {\thesubsection}
    {1em}
    {}

\titleformat{\subsubsection}
    {\normalsize\bfseries\itshape}
    {\thesubsubsection}
    {1em}
    {}

\titlespacing*{\section}{0pt}{2.5ex plus 1ex minus .2ex}{1.5ex plus .2ex}
\titlespacing*{\subsection}{0pt}{2ex plus 1ex minus .2ex}{1ex plus .2ex}
\titlespacing*{\subsubsection}{0pt}{1.5ex plus 1ex minus .2ex}{0.5ex plus .2ex}

% =============================================================================
% HYPERLINKS (load near end)
% =============================================================================
\usepackage[
    colorlinks=true,
    linkcolor=thesisblue,
    citecolor=citegreen,
    urlcolor=linkblue,
    bookmarks=true,
    bookmarksnumbered=true,
    pdfstartview=FitH,
    pdftitle={Graph-Structured Adaptive Memory: Evolving Contexts for Agentic, Self-Improving Systems},
    pdfauthor={Machungo Erick},
    pdfsubject={TDK Thesis - University of Debrecen},
    pdfkeywords={LLM, knowledge graphs, context engineering, GSAM, financial NLP}
]{hyperref}

% =============================================================================
% TYPOGRAPHY REFINEMENTS
% =============================================================================
\widowpenalty=10000
\clubpenalty=10000

% Paragraph formatting
\setlength{\parindent}{1.25cm}
\setlength{\parskip}{0pt}

% =============================================================================
% CUSTOM COMMANDS
% =============================================================================
\newcommand{\imgplaceholder}[1]{%
    \textcolor{gray}{\small\textit{[Placeholder: #1]}}%
}

% Convenience commands
\newcommand{\eg}{e.g.,\xspace}
\newcommand{\ie}{i.e.,\xspace}
\newcommand{\cf}{cf.\xspace}
\newcommand{\etal}{et al.\xspace}

% =============================================================================
% HEADERS AND FOOTERS
% =============================================================================
\usepackage{fancyhdr}

\pagestyle{fancy}
\fancyhf{}
\fancyhead[L]{\small\nouppercase{\leftmark}}
\fancyhead[R]{\small\thepage}
\renewcommand{\headrulewidth}{0.4pt}
\renewcommand{\footrulewidth}{0pt}

% Plain style for chapter/section first pages
\fancypagestyle{plain}{
    \fancyhf{}
    \fancyfoot[C]{\small\thepage}
    \renewcommand{\headrulewidth}{0pt}
}

% =============================================================================
\begin{document}

% =============================================================================
% TITLE PAGE
% =============================================================================
\begin{titlepage}
\thispagestyle{empty}

\begin{center}

% University header
{\Large\scshape University of Debrecen}\\[0.3cm]
{\large\scshape Faculty of Informatics}\\[0.2cm]
{\scshape Department of Computer Science}

\vspace{0.5cm}

% University logo (after department)
\includegraphics[width=3.5cm]{university_logo.png}

\vspace{0.8cm}

% Title (no blue lines)
{\fontsize{16}{20}\selectfont\bfseries
Graph-Structured Adaptive Memory:\\[0.3cm]
Evolving Contexts for Agentic,\\[0.3cm]
Self-Improving Systems
}

\end{center}

\vspace{0.8cm}

% Author and Supervisor
\begin{minipage}[t]{0.48\textwidth}
\begin{flushleft}
\textbf{Supervisor:}\\[0.2cm]
Dr.~Bogacsovics Gerg\H{o}\\
\textit{Professor}\\
Department of Computer Science
\end{flushleft}
\end{minipage}
\hfill
\begin{minipage}[t]{0.48\textwidth}
\begin{flushright}
\textbf{Author:}\\[0.2cm]
Machungo Erick\\
\textit{Computer Science BSc}\\
Faculty of Informatics
\end{flushright}
\end{minipage}

\vspace{0.8cm}

% Consent
\begin{center}
\textit{``I consent to its submission for review''}

\vspace{0.6cm}

\rule{5cm}{0.4pt}\\[0.2cm]
{\small Supervisor's signature}

\vspace{0.8cm}

% Place and year
{\large Debrecen, 2026}

\end{center}

\end{titlepage}

% =============================================================================
% TABLE OF CONTENTS
% =============================================================================
\newpage
\thispagestyle{plain}
\tableofcontents
\newpage

% =============================================================================
% INTRODUCTION
% =============================================================================
\section{Introduction}
\label{sec:introduction}

Large language models (LLMs) are now essential tools in almost every area of computing, which includes software engineering, scientific research, financial analysis, and legal reasoning \citep{brown2020language, openai2023gpt4}. Even though LLMs are very good at a lot of things, they have difficulties with specific domains and areas that require accurate, and detailed knowledge. These areas have knowledge that is constantly changing and needs to be up-to-date. Large language models are not as good, in these domains. They struggle with structured information. This constitutes a constraint of language models. A model trained on broad internet data does not inherently understand the hierarchical relationships within the XBRL (eXtensible Business Reporting Language) taxonomy used in financial reporting, nor can it reason about how an error in identifying one financial entity propagates through downstream formulas \citep{loukas2022finer}.

Two broad strategies have emerged to address this gap. \textit{Fine-tuning} updates the model's weights on domain-specific data \citep{wang2025finlora}, but is expensive, non-interpretable, and difficult to update incrementally. \textit{Context adaptation} (also called context engineering) instead modifies the \textit{inputs} to the model---injecting instructions, strategies, examples, or evidence into the prompt without altering weights \citep{ace2025}. Context adaptation is interpretable (the injected knowledge is readable text), cheap to update (no retraining), and shareable across models.

Recent context adaptation frameworks have shown impressive results. Dynamic Cheatsheet \citep{suzgun2025dynamic} introduced the idea of an evolving memory of strategies accumulated during test-time inference. ACE (Agentic Context Engineering) \citep{ace2025} addressed two critical failure modes of earlier approaches---\textit{brevity bias} (where optimization collapses toward overly concise summaries that lose domain detail) and \textit{context collapse} (where iterative LLM-driven rewriting gradually degrades accumulated knowledge)---through a multi-agent architecture with incremental delta updates. ACE achieved average gains of 8.6\% on financial benchmarks and 10.6\% on agent tasks, while reducing adaptation latency by 86.9\% compared to prior methods \citep{ace2025}.

However, both Dynamic Cheatsheet and ACE share a fundamental architectural limitation: they store accumulated knowledge as \textit{flat collections of bullet points}. Each bullet is an independent unit with metadata (helpful and harmful counters), but with no explicit connections to other bullets. This design discards a signal that is critically important in specialized domains: the \textit{relationships between knowledge elements}.

\subsection{Research Objectives}

In this thesis, I introduce \textbf{GSAM (Graph-Structured Adaptive Memory)}, a framework that replaces flat bullet-point storage with an ontology-grounded knowledge graph. The GSAM framework builds upon the ACE framework in the following ways. The framework extends ACE's Generator--Reflector--Curator architecture with two new components: a \textbf{Graph Constructor} and a \textbf{Graph Retriever}. The Graph Constructor mainly builds and maintains the knowledge graph from learning experiences. The Graph Retriever's main task is to extract relevant subgraphs for new tasks through a combination of embedding similarity and structural graph traversal.

The research targets three hypotheses:

\begin{description}[style=nextline, leftmargin=1.5cm]
\item[\textbf{H1:}] Graph-structured memory with ontology grounding improves cross-concept knowledge transfer compared to flat storage.
\item[\textbf{H2:}] Explicit failure cascade modelling reduces repeated and downstream errors.
\item[\textbf{H3:}] Ontology-aware retrieval achieves higher precision than embedding-only retrieval.
\end{description}

I evaluate GSAM on two main financial analysis benchmarks: FiNER \citep{loukas2022finer} (139-type entity recognition) and Formula \citep{wang2025finlora} (numerical reasoning over financial filings). Additionally, I introduce \textbf{FiNER-Transfer}, a new benchmark that is specifically designed to measure cross-concept transfer within the XBRL taxonomy. The choice to focus exclusively on domain-specific financial tasks, rather than general agent benchmarks such as AppWorld \citep{trivedi2024appworld}, is deliberate. Financial analysis provides formal ontological structure (the XBRL taxonomy), natural dependency graphs (formulas depend on entity extraction), and compositional error cascades---all of which are naturally and easily captured by graphs but are completely invisible when it comes to flat storage used in both DC and ACE.

\subsection{Thesis Structure}

The rest of this thesis is organized as follows. Section~\ref{sec:background} reviews the relevant literature of context adaptation, knowledge graphs, and financial natural language processing (NLP). Section~\ref{sec:gsam} presents the GSAM framework including the knowledge graph schema, ontology layer, failure cascade model, and retrieval mechanism. Section~\ref{sec:finer-transfer} describes how to construct the FiNER-Transfer benchmark. Section~\ref{sec:experiments} shows the details of the experimental design and its results. Section~\ref{sec:cost} analyzes the latency and cost implications of graph-structured memory. Section~\ref{sec:discussion} discusses the findings, problems, and next steps. Section~\ref{sec:summary} reveals the summary of the key contributions.

% =============================================================================
% BACKGROUND AND RELATED WORK
% =============================================================================

\section{Background and Related Work}
\label{sec:background}

\subsection{Large Language Models and Context Windows}

The transformer architecture \citep{vaswani2017attention} created a new way to process the natural language. It begins by using self-attention mechanisms to process input sequences instead of recurrence. When this architecture was scaled up to billions of parameters, it led to models like GPT-3 \citep{brown2020language}, GPT-4 \citep{openai2023gpt4}, LLaMA \citep{touvron2023llama}, and DeepSeek \citep{deepseek2024}, which exhibit strong few-shot learning capabilities on a wide range of tasks.

The key insight underlying few-shot learning is that the model's behavior can be substantially altered by what appears in the \textit{context window} which is the input provided alongside the query \citep{brown2020language}. This observation motivated a shift from fine-tuning (modifying weights) toward context engineering (modifying inputs) as a more flexible approach to domain adaptation. Retrieval-augmented generation (RAG) \citep{lewis2020retrieval} gets relevant documents to add to the context. Chain-of-thought prompting \citep{wei2022chain} creates step-by-step reasoning by putting the context together in the correct way. Decomposed prompting \citep{khot2022decomposed} divides complicated tasks into smaller sub-tasks. These methods show that the context is just as important as how the model was trained.

\subsection{From Dynamic Cheatsheet to ACE}

The progress we have seen in evolution from the use of static prompts to the application of evolving memory systems represents the current frontier in context engineering. \textbf{Dynamic Cheatsheet} \citep{suzgun2025dynamic} introduced the concept of a ``cheatsheet''. A cheat sheet can be thought of as an external memory of reusable strategies that accumulates during test-time inference. As the model encounters new tasks, it reflects on its performance and adds all successful strategies it has encountered to the cheatsheet. This simple idea produced substantial and notable improvements, but still suffered from two problems. First, \textit{brevity bias}: the LLM tasked with maintaining the cheatsheet tended to compress entries into progressively shorter summaries, which results in losing important and crucial domain-specific detail. Second, \textit{context collapse}: this happens through the iterative rewriting by the same LLM, which gradually degrades the quality of accumulated knowledge.

\textbf{ACE (Agentic Context Engineering)} \citep{ace2025} addressed both problems that Dynamic Cheatsheet (DC) was facing through architectural innovation. Rather than having a single model manage the entire memory, ACE took a different approach by separating responsibilities across three specialized agents. The three agents are a Generator agent, a Reflector agent and lastly a Curator agent:

\begin{itemize}[leftmargin=*]
\item The \textbf{Generator agent} produces reasoning trajectories for new queries, using the current knowledge context to inform its approach.
\item The \textbf{Reflector agent} analyzes the trajectory from the generator agent against feedback (correct/incorrect), extracting concrete insights about what worked and what failed.
\item The \textbf{Curator agent} finally converts the Reflector agent's insights into structured delta updates. Delta updates include specific additions, modifications, or removals that are applied to the knowledge store using deterministic merging with de-duplication and pruning.
\end{itemize}

This clever approach of separation prevents context collapse because the Curator agent never rewrites the entire context. The Curator agent only applies localized edits. ACE also introduces a ``grow-and-refine'' principle. This simply means that the context grows by accumulating detailed entries, and is refined by merging near-duplicates and pruning low-utility items. On the financial benchmarks (FiNER and Formula), ACE achieved an average improvement of 8.6\%. This improvement is over strong baselines, including in-context learning and GEPA \citep{agrawal2025gepa}. ACE achieved this improvement while reducing adaptation latency by 91.5\% compared to Dynamic Cheatsheet \citep{ace2025}.

\subsection{The Flat Storage Bottleneck}
\label{sec:flat-bottleneck}

Despite ACE's effectiveness, its knowledge store is fundamentally a flat list of bullet points. Each bullet is an independent text entry with metadata---a helpful count, a harmful count, and a section label---but no explicit connections to other bullets. I identify three limitations inherent to this representation:

\textbf{Independent Retrieval.} When the Generator faces a new task, relevant bullets are retrieved via embedding cosine similarity. Each bullet competes independently for the top-$k$ retrieval slots. This means that a strategy about extracting \texttt{NetIncome} and a related anti-pattern about confusing \texttt{NetIncome} with \texttt{ComprehensiveIncome} are retrieved independently---there is no mechanism to ensure that retrieving one also retrieves the other, even though they are about the same concept.

\textbf{Flat Namespace.} All bullets share a single semantic space. Domain concepts, strategies, troubleshooting tips, and formula dependencies compete for the same retrieval slots. As the number of bullets grows, the embedding space becomes crowded and retrieval precision degrades.

\textbf{No Structural Priors.} Financial analysis is organized around formal structures---the XBRL taxonomy defines a hierarchy of 139+ entity types, and formulas define computational dependencies between entities. Flat storage cannot leverage these structures for knowledge organization or transfer.

\subsection{Knowledge Graphs for Language Models}

Knowledge graphs (KGs) represent structured information as typed entities and relations \citep{hogan2021knowledge}. They have been successfully applied to enhance LLM capabilities in several settings. QA-GNN \citep{yasunaga2021qagnn} fuses language models with knowledge graphs for question answering by jointly reasoning over text and graph structure. KAPING \citep{baek2023kaping} retrieves KG subgraphs to augment prompts in a zero-shot setting. KG-Agent \citep{jiang2024kgagent} uses knowledge graphs to guide multi-step reasoning in autonomous agents.

However, all of these approaches use \textit{pre-existing} knowledge graphs---typically large-scale resources such as ConceptNet \citep{speer2017conceptnet} or Wikidata. GSAM is fundamentally different in that it builds a knowledge graph dynamically from the LLM's own learning process. It is based on a domain ontology instead of consuming a static external knowledge graph.

\subsection{Agent Memory Systems}

A growing body of work explores memory architectures for LLM agents. AgentFly \citep{zhou2025agentfly} enables continuous memory evolution across tasks. AWM \citep{wang2024awm} induces reusable workflows from agent trajectories. A-MEM \citep{xu2025amem} organizes memory using structured attributes inspired by the Zettelkasten method. MemPrompt \citep{madaan2022memory} stores success and failure examples for retrieval. ExpeL \citep{zhao2024expel} learns task-specific expertise from accumulated experience. Generative agents \citep{park2023generative} use a memory stream with reflection to simulate human-like behavior.

These systems all store memories as individual entries (experiences, workflows, or observations). GSAM's contribution is to model the \textit{relationships between} memory elements explicitly, and to ground those relationships in a formal domain ontology.

\subsection{Financial NLP and the XBRL Standard}
\label{sec:xbrl-background}

Financial natural language processing has received growing attention as a testbed for domain-specific reasoning. \textbf{FiNER} \citep{loukas2022finer} introduced the task of financial numeric entity recognition: given a sentence from a financial filing, identify which of 139 XBRL entity types each numeric value represents. The 139 types are organized in a hierarchical taxonomy defined by the XBRL standard---for example, \texttt{Assets} $\rightarrow$ \texttt{CurrentAssets} $\rightarrow$ \texttt{CashAndCashEquivalents}.

The XBRL (eXtensible Business Reporting Language) standard is an international framework for digital business reporting, maintained by XBRL International. It provides a formal taxonomy of financial concepts with explicit hierarchical relationships, making it an ideal testbed for evaluating whether graph-structured memory can exploit domain structure. The \textbf{Formula} benchmark \citep{wang2025finlora} complements FiNER by testing numerical reasoning over XBRL filings by calculating financial ratios and aggregates that are dependent upon accurately extracted entities.

% =============================================================================
\section{The GSAM Framework}
\label{sec:gsam}
% =============================================================================

GSAM enhances the ACE architecture and design by replacing the flat bullet store with a graph-structured memory module. The three ACE agents (Generator, Reflector, Curator) are maintained, and two new components are added. The first component is the \textbf{Graph Constructor}, which converts Curator output into graph operations, while the second component is the \textbf{Graph Retriever}, which extracts relevant subgraphs for new tasks. Figure~\ref{fig:architecture} illustrates the full architecture.

\begin{figure}[H]
\centering
\includegraphics[width=\textwidth]{gsam_framework.png}
\caption{The GSAM architecture. The Generator, Reflector, and Curator were inherited from ACE. GSAM includes a Graph Constructor (which transforms Curator deltas into graph operations) and a Graph Retriever (which extracts relevant subgraphs for each new query). The Knowledge Graph $G$ replaces ACE's flat bullet store. Arrows show data flow: the query enters the Generator with retrieved context, the trajectory goes through the Reflector and Curator, and the resulting deltas are converted to graph operations by the Graph Constructor.}
\label{fig:architecture}
\end{figure}

\subsection{Problem Formulation}

Let $\mathcal{T} = \{t_1, \ldots, t_N\}$ be a sequence of domain-specific tasks. Each task $t_i = (c_i, q_i, y_i)$ consists of a context $c_i$ (e.g., a financial document excerpt), a question $q_i$, and a ground-truth answer $y_i$. The system produces a prediction $\hat{y}_i$ and receives feedback $f_i$ (correct/incorrect, with optional reasoning trace).

\begin{definition}[Flat Storage]
In ACE, knowledge is a set $\mathcal{B} = \{b_1, \ldots, b_M\}$ of independent bullets. Each bullet $b_j = (\text{content}_j, h_j, d_j, s_j)$ contains textual content, a helpful count $h_j$, a harmful count $d_j$, and a section label $s_j$. After each task, the Curator produces delta bullets $\Delta\mathcal{B}_i$ which are merged: $\mathcal{B} \leftarrow \mathcal{B} \cup \Delta\mathcal{B}_i$.
\end{definition}

\begin{definition}[Graph-Structured Memory]
In GSAM, knowledge is a typed, attributed directed graph $G = (V, E, \tau_V, \tau_E, \alpha)$ where $V$ is a set of nodes, $E \subseteq V \times V$ is a set of directed edges, $\tau_V : V \rightarrow \mathcal{T}_V$ assigns a type to each node, $\tau_E : E \rightarrow \mathcal{R}$ assigns a type to each edge, and $\alpha : V \cup E \rightarrow \mathcal{A}$ maps elements to attribute dictionaries containing natural language content, counters, and confidence scores.
\end{definition}

The node type set for financial analysis is $\mathcal{T}_V = \{\textsc{Strategy}, \textsc{AntiPattern}, \textsc{Concept}, \textsc{Formula}, \textsc{Confusion}\}$, and the relation set is:
\begin{equation}
\mathcal{R} = \left\{
\begin{array}{l}
\texttt{is\_a},\ \texttt{part\_of},\ \texttt{depends\_on},\ \texttt{applies\_to},\\
\texttt{fails\_for},\ \texttt{fixes},\ \texttt{confused\_with},\ \texttt{conflicts\_with}
\end{array}
\right\}
\label{eq:relations}
\end{equation}

Figure~\ref{fig:flat-vs-graph} illustrates the difference between ACE's flat storage and GSAM's graph-structured memory for the same set of financial knowledge.

\begin{figure}[H]
\centering
\includegraphics[width=\textwidth]{figure2}
\caption{Comparison of ACE's flat bullet storage (left) and GSAM's knowledge graph (right) for the same financial analysis knowledge. In flat storage, each bullet is an independent entry retrieved by embedding similarity. In GSAM, strategies and anti-patterns are connected to the domain concepts they address, enabling structural retrieval and cross-concept transfer.}
\label{fig:flat-vs-graph}
\end{figure}

\subsection{Knowledge Graph Schema}
\label{sec:schema}

Table~\ref{tab:schema} defines the complete node and edge types in the GSAM schema. A key design decision is the separation between \textit{domain-structural} edges (\texttt{is\_a}, \texttt{part\_of}, \texttt{depends\_on}), which encode the formal ontology and can be initialized before any tasks are executed, and \textit{experiential} edges (\texttt{applies\_to}, \texttt{fails\_for}, \texttt{fixes}, \texttt{confused\_with}, \texttt{conflicts\_with}), which are learned from execution feedback during adaptation.

\begin{table}[H]
\centering
\caption{GSAM knowledge graph schema for financial analysis.}
\label{tab:schema}
\small
\begin{tabular}{@{}llp{7cm}@{}}
\toprule
\textbf{Component} & \textbf{Type} & \textbf{Description and Attributes} \\
\midrule
\multicolumn{3}{l}{\textit{Node Types}} \\
\midrule
\textsc{Strategy} & Node & Reusable approach or pattern. Attributes: content, helpful count, harmful count, confidence score. \\
\textsc{AntiPattern} & Node & Failed approach to avoid. Attributes: content, failure count, root cause description, severity level. \\
\textsc{Concept} & Node & Domain concept (e.g., XBRL entity type). Attributes: name, definition, taxonomy path. \\
\textsc{Formula} & Node & Computational rule. Attributes: expression, input entities, output entity. \\
\textsc{Confusion} & Node & Documented confusion between entities. Attributes: entity pair, distinguishing criteria. \\
\midrule
\multicolumn{3}{l}{\textit{Edge Types (Domain-Structural)}} \\
\midrule
\texttt{is\_a} & Edge & Taxonomic: child $\rightarrow$ parent in XBRL hierarchy. \\
\texttt{part\_of} & Edge & Compositional: component $\rightarrow$ aggregate concept. \\
\texttt{depends\_on} & Edge & Computational: formula $\rightarrow$ required input entity. \\
\midrule
\multicolumn{3}{l}{\textit{Edge Types (Experiential)}} \\
\midrule
\texttt{applies\_to} & Edge & Strategy $\rightarrow$ Concept it addresses. \\
\texttt{fails\_for} & Edge & AntiPattern $\rightarrow$ Concept where it manifests. \\
\texttt{fixes} & Edge & Strategy $\rightarrow$ AntiPattern it resolves. \\
\texttt{confused\_with} & Edge & Concept $\leftrightarrow$ Concept (bidirectional confusion). \\
\texttt{conflicts\_with} & Edge & Strategy $\leftrightarrow$ Strategy (mutually exclusive approaches). \\
\bottomrule
\end{tabular}
\end{table}

\subsection{Domain Ontology Layer}
\label{sec:ontology}

A distinguishing feature of GSAM is its use of a formal domain ontology as a structural backbone for the knowledge graph. Before any adaptation begins, the graph is pre-populated with concept nodes and taxonomic edges derived from the XBRL taxonomy.

\textbf{Ontology Initialization.} For FiNER, this creates 139 \textsc{Concept} nodes organized hierarchically. For each parent--child relationship in the XBRL taxonomy, we add a directed \texttt{is\_a} edge:
\begin{equation}
\forall\, (c_{\text{child}},\ c_{\text{parent}}) \in \text{XBRL\_Taxonomy}: \quad (c_{\text{child}}, \texttt{is\_a}, c_{\text{parent}}) \in E
\end{equation}

For Formula tasks, we additionally extract \textsc{Formula} nodes from the training data and link them to their input concepts via \texttt{depends\_on} edges. Figure~\ref{fig:xbrl} shows the hierarchical structure of the XBRL taxonomy used in this work.

\begin{figure}[H]
\centering
\includegraphics[width=0.85\textwidth]{figure3}
\caption{Excerpt of the XBRL taxonomy hierarchy used to initialize the Domain Ontology Layer. The taxonomy defines parent--child relationships among financial concepts (e.g., \texttt{Assets} $\rightarrow$ \texttt{CurrentAssets} $\rightarrow$ \texttt{CashAndCashEquivalents}). These relationships become \texttt{is\_a} edges in the knowledge graph.}
\label{fig:xbrl}
\end{figure}

\textbf{Ontology-Guided Transfer.} This is the taxonomic backbone that enables us to have some sort of grounds to form a structural transfer, which is otherwise unavailable in flat storage. Here is an illustration, if a strategy $s$ has been learned for concept $c_1$, and concepts $c_1$ and $c_2$ share a common parent $p$ in the taxonomy (i.e., they are siblings), then $s$ is a \textit{candidate} for transfer to tasks involving $c_2$:
\begin{equation}
\text{transfer\_candidates}(c_2) = \{s : (s, \texttt{applies\_to}, c_1) \in E \land \text{lca}(c_1, c_2) \leq d_{\max}\}
\label{eq:transfer}
\end{equation}
where $\text{lca}(c_1, c_2)$ denotes the depth of the lowest common ancestor of $c_1$ and $c_2$ in the taxonomy, and $d_{\max}$ is a threshold controlling transfer breadth. Strategies that are transferred in this way are flagged as \textit{tentative} until they are validated by the execution feedback on the concept being targeted.

\subsection{Failure Cascade Model}
\label{sec:failure}

In flat storage, each failure is recorded as an independent bullet. GSAM represents failures as nodes within a connected structure that captures the reasons for failure and their consequences.

\textbf{Anti-Pattern Nodes.} When the Reflector identifies a failure, the Graph Constructor creates an \textsc{AntiPattern} node with rich attributes. They are the failure description, a root cause analysis, a severity level, and an occurrence count.

\begin{lstlisting}[caption={Example AntiPattern node attributes.},label={lst:antipattern}]
AntiPattern: "Confusing NetIncome with ComprehensiveIncome"
  root_cause: "Both appear in income statements;
    ComprehensiveIncome includes unrealized gains/losses"
  severity: high
  frequency: 3
\end{lstlisting}

\textbf{Failure Propagation Edges.} The Graph Constructor connects anti-patterns to all affected downstream parts. Critically, this includes not only the directly confused concept but also every formula that depends on that concept:
\begin{equation}
\text{affected}(a) = \{f : (a, \texttt{fails\_for}, c) \in E \land (f, \texttt{depends\_on}, c) \in E\}
\end{equation}

\textbf{Cascade Detection.} Before performing a new task with a formula $f$, GSAM checks for upstream anti-patterns by examining dependency edges:
\begin{equation}
\text{risks}(f) = \{a : \exists\, c\ \text{s.t.}\ (f, \texttt{depends\_on}, c) \in E \land (a, \texttt{fails\_for}, c) \in E\}
\end{equation}
These are shown to the Generator as explicit warnings, which enable to avoid the errors before they happen. Through this method, GSAM achieves hypothesis H2: by showing failure cascades which lead the system to be able to prevent errors before they happen.

Figure~\ref{fig:knowledge-graph} shows an example of a knowledge subgraph illustrating how concept nodes, strategies, anti-patterns, and formulas are interconnected through typed edges.

\begin{figure}[H]
\centering
\includegraphics[width=0.9\textwidth]{figure4}
\caption{Example knowledge subgraph centered on the concept \texttt{NetIncome}. The subgraph shows a \textsc{Strategy} node connected via \texttt{applies\_to}, an \textsc{AntiPattern} node connected via \texttt{fails\_for}, a \textsc{Formula} node (\texttt{EPS}) connected via \texttt{depends\_on}, and a \texttt{confused\_with} edge to \texttt{ComprehensiveIncome}. The \texttt{fixes} edge connects the strategy to the anti-pattern it resolves.}
\label{fig:knowledge-graph}
\end{figure}

\subsection{Graph Construction Pipeline}
\label{sec:construction}

GSAM's graph construction extends ACE's Curator with a Graph Constructor module that converts natural-language delta insights into graph operations. Algorithm~\ref{alg:construction} presents the complete pipeline.

\begin{algorithm}[H]
\caption{GSAM Graph Construction Pipeline}
\label{alg:construction}
\begin{algorithmic}[1]
\REQUIRE Task $t_i$, trajectory $\tau_i$, feedback $f_i$, current graph $G$
\STATE $\text{insights} \leftarrow \text{Reflector}(\tau_i, f_i, G)$
\STATE $\text{deltas} \leftarrow \text{Curator}(\text{insights}, G)$
\STATE $\text{operations} \leftarrow \text{GraphConstructor}(\text{deltas})$
\FOR{each operation $\text{op}$ in operations}
    \IF{op.type = \texttt{ADD\_NODE}}
        \STATE $v_{\text{new}} \leftarrow \text{CreateNode}(\text{op.node\_type}, \text{op.content})$
        \IF{$\exists\, v \in V : \text{sim}(v, v_{\text{new}}) > \theta_{\text{merge}} \land \tau_V(v) = \tau_V(v_{\text{new}})$}
            \STATE $\text{MergeNodes}(v, v_{\text{new}})$ \COMMENT{De-duplication via semantic similarity}
        \ELSE
            \STATE $V \leftarrow V \cup \{v_{\text{new}}\}$
        \ENDIF
    \ELSIF{op.type = \texttt{ADD\_EDGE}}
        \STATE $E \leftarrow E \cup \{(\text{op.source}, \text{op.relation}, \text{op.target})\}$
    \ELSIF{op.type = \texttt{UPDATE\_ATTR}}
        \STATE Update $\alpha(\text{op.node})$ with op.updates \COMMENT{Increment counters}
    \ENDIF
\ENDFOR
\STATE $G \leftarrow \text{Prune}(G, \theta_{\text{prune}})$ \COMMENT{Remove low-utility nodes}
\RETURN $G$
\end{algorithmic}
\end{algorithm}

The Graph Constructor uses a structured LLM prompt with few-shot examples from the financial domain (see Appendix~\ref{app:prompt}) to extract entities and relationships from the Curator's natural-language deltas. Extracted nodes are de-duplicated against existing graph nodes using cosine similarity with threshold $\theta_{\text{merge}} = 0.9$, requiring matching node types. When a new node is sufficiently similar to an existing node of the same type, the two are merged: their content is combined and their counters are aggregated.

\textbf{Ontology Anchoring.} When a \textsc{Strategy} or \textsc{AntiPattern} mentions a domain concept, the Graph Constructor creates \texttt{applies\_to} or \texttt{fails\_for} edges linking to the corresponding \textsc{Concept} node in the ontology. This anchoring connects experiential knowledge to the structural backbone.

\textbf{Pruning.} Nodes with degree less than 2 and confidence below $\theta_{\text{prune}}$ for more than $K$ epochs are removed along with their incident edges.

\subsection{Multi-Epoch Graph Refinement}
\label{sec:multiepoch}

ACE supports multi-epoch adaptation, where the same training queries are revisited across multiple passes. In ACE's ablation study, removing multi-epoch reduced accuracy substantially \citep{ace2025}, confirming that second-pass encounters refine incomplete strategies.

GSAM inherits multi-epoch support but introduces \textit{graph-specific refinement operations} that go beyond what is possible with flat bullets. In flat storage, a second encounter with query $q$ can only add new bullets or increment counters. In GSAM, a second encounter triggers three additional operations:

\textbf{Edge Discovery.} On the first encounter, the Graph Constructor may create a \textsc{Strategy} node linked to the most salient \textsc{Concept}. On a second encounter with a related query, the system may discover that the strategy also applies to additional concepts, creating new \texttt{applies\_to} edges that were invisible on the first pass.

\textbf{Edge Weight Reinforcement.} When a strategy is reused successfully, the weight on its \texttt{applies\_to} edge is incremented. When it fails in a specific context, the corresponding \texttt{fails\_for} edge is created or strengthened:
\begin{equation}
w(s, \texttt{applies\_to}, c) \leftarrow w(s, \texttt{applies\_to}, c) + \mathbb{1}[\text{helpful in context } c]
\end{equation}

\textbf{Cross-Node Consolidation.} After each epoch, a consolidation pass identifies nodes that have acquired similar neighborhoods (connected to the same concepts via the same edge types) and merges them, producing more general strategies than would emerge from any single encounter.

The net and overall effect is that multi-epoch adaptation in the Graph Structured Adaptive Memory (GSAM) refines both the \textit{content} of nodes (as in ACE) and the \textit{structure} of the graph. This results in the discovery of relationships, validation of connections, and merging redundant knowledge.

\subsection{Ontology-Aware Retrieval}
\label{sec:retrieval}

Let us see how GSAM uses ontology-grounded reasoning to retrieve knowledge. Let us say that you are given a new task $t_i = (c_i, q_i, y_i)$, GSAM retrieves a relevant subgraph $G_i$ through a three-stage process:

\textbf{Stage 1: Concept Identification.} In this stage, the system extracts domain concepts mentioned in the query $q_i$ and context $c_i$ using both keyword matching against concept names and embedding similarity:
\begin{equation}
\mathcal{C}_i = \{v \in V_{\textsc{Concept}} : \text{sim}(q_i \oplus c_i,\ v) > \theta_{\text{concept}}\}
\end{equation}

\textbf{Stage 2: Graph Traversal.} In this stage, from the matched concepts $\mathcal{C}_i$, the system performs breadth-first search (BFS) up to depth $d = 2$, collecting strategies, anti-patterns, formulas, and dependencies:
\begin{align}
\mathcal{S}_i &= \{v : (v, \texttt{applies\_to}, c) \in E,\ c \in \mathcal{C}_i\} & \text{(strategies)} \\
\mathcal{A}_i &= \{v : (v, \texttt{fails\_for}, c) \in E,\ c \in \mathcal{C}_i\} & \text{(anti-patterns)} \\
\mathcal{D}_i &= \{v : (f, \texttt{depends\_on}, v) \in E,\ f \in \mathcal{S}_i\} & \text{(dependencies)}
\end{align}

\textbf{Stage 3: Taxonomic Expansion.} In this stage, for each matched concept, the system includes strategies from taxonomically related (sibling) concepts:
\begin{equation}
\mathcal{S}_i^+ = \mathcal{S}_i \cup \{s : (s, \texttt{applies\_to}, c') \in E \land (c', \texttt{is\_a}, \text{parent}(c)) \in E\}
\end{equation}
These transfer candidates are included with a reduced confidence weight and flagged as tentative.

\textbf{Serialization.} Last but not least, the retrieved subgraph is then serialized into structured natural language that the Generator agent can consume. The serialization process organizes information by concept, grouping each concept's strategies, anti-patterns, and formula dependencies, followed by a separate section for transferred strategies from sibling concepts. This structure gives the Generator agent a well organized and concept-centred view of all relevant knowledge, rather than a flat list of bullets sorted only by embedding similarity.

% =============================================================================
\section{FiNER-Transfer: A Cross-Concept Transfer Benchmark}
\label{sec:finer-transfer}
% =============================================================================

In order to directly test hypothesis H1 (cross-concept transfer), I introduce \textbf{FiNER-Transfer}, a controlled benchmark that is derived from the XBRL taxonomy structure. FiNER-Transfer was made because all the existing benchmarks assess the total and overall accuracy of a task, but they do not distinguish whether the knowledge acquired for one concept is or can apply to any related concepts.

\subsection{Benchmark Construction}

There are four steps to building FiNER-Transfer:

\textbf{Step 1: Taxonomy Clustering.} The 139 FiNER entity types are put into grouped based on their position in the XBRL hierarchy. Sibling concepts that share a direct parent form concept pairs. I find 42 such pairs in the taxonomy.

\textbf{Step 2: Task Pair Selection.} From the FiNER dataset, I select examples that involve each concept in each pair. For each pair $(c_1, c_2)$, there is a source set with examples of $c_1$ and a target set with examples of $c_2$.

\textbf{Step 3: Transfer Protocol.} The system is first trained on source examples (learning strategies for $c_1$), then is tested on target examples ($c_2$) without additional training. The transfer gain is:
\begin{equation}
\Delta_{\text{transfer}}(c_1 \rightarrow c_2) = \text{Acc}(c_2 \mid \text{adapted on } c_1) - \text{Acc}(c_2 \mid \text{no adaptation})
\end{equation}

\textbf{Step 4: Controls.} Distant pairs which are ideas from different subtrees of the taxonomy (e.g., a pair of Balance Sheet concept and an Income Statement concept), act as negative controls where minimal transfer is expected.

\subsection{Evaluation Metrics}

The newly introduced benchmark, FiNER-Transfer uses four main metrics. The metrics are \textbf{Near-Transfer Rate}, \textbf{Far-Transfer Rate}, \textbf{Transfer Precision} and lastly \textbf{Negative Transfer Rate}. Let us start with the Near-Transfer Rate metric: this constitutes the fraction of sibling pairs with positive transfer. The second metric, Far-Transfer Rate, constitutes the fraction of distant pairs with positive transfer. The third metric, Transfer Precision, constitutes the average $\Delta_{\text{transfer}}$ among positive transfers. Lastly, the Negative Transfer Rate metric constitutes the fraction of pairs where adaptation \textit{hurts} performance on the target concept.

% =============================================================================
\section{Experimental Design and Results}
\label{sec:experiments}
% =============================================================================

\subsection{Datasets}

Let me start by listing and expanding on all the datasets used when running the experiments.

\textbf{FiNER} \citep{loukas2022finer}: Financial entity recognition with 139 XBRL entity types. Briefly, this is a dataset and a framework designed to automate eXtensive Business Reporting Language (XBRL) tagging in SEC (Securities and Exchange Commission) filings. The dataset consists of up to 1.1 million sentences that are carefully annotated with the 139 types of numeric tags, utilizing context to identify and label financial figures for improved automated reporting. For our experiment specifically, the FiNER dataset consists of up to 1,000 training and a total of 300 test examples. The main metric measured and evaluated is the exact-match accuracy.

\textbf{Formula} \citep{wang2025finlora}: Numerical reasoning over XBRL filings. In summary, this is a framework made specifically for numerical reasoning over structured XBRL financial filings. The framework focuses on selecting relevant facts, constructing financial formulas, and computing results. This approach enhances LLM's precision in financial reporting by utilizing agentic frameworks to navigate the taxonomy-driven data and address complex calculations. For our experiment specifically, the Formula dataset contains a total of 800 training examples and a total of 200 test examples. The metric being evaluated here is also exact-match accuracy.

\textbf{FiNER-Transfer} (introduced in this thesis): 42 sibling concept pairs derived from the XBRL taxonomy, plus 42 distant pairs as controls.

\subsection{Baselines and Methods}

Every method in these experiments runs on \textbf{DeepSeek-R1-Distill-Llama-70B} \citep{deepseek2024} deployed via a self-hosted vLLM server on Modal cloud infrastructure. Because the ACE paper \citep{ace2025} used the larger DeepSeek-V3 model, absolute accuracy values in our experiments may differ from those originally reported; all comparisons between ACE and GSAM are however conducted on identical infrastructure with the same model, making the relative differences meaningful. The baselines are:

\begin{itemize}[leftmargin=*]
\item \textbf{Base LLM}: Direct evaluation that doesn't use any context engineering.
\item \textbf{ICL} \citep{brown2020language}: In-context learning uses training examples chosen at random.
\item \textbf{GEPA} \citep{agrawal2025gepa}: Genetic Pareto prompt optimization uses evolutionary strategies to find the best prompt variants.
\item \textbf{ACE} \citep{ace2025}: Agentic context engineering with flat bullet storage (primary baseline) which is used to compare to GSAM.
\item \textbf{GSAM} (this thesis): Graph-structured adaptive memory with ontology grounding.
\end{itemize}

\subsection{Experimental Protocol}

I follow ACE's evaluation protocol to ensure fair comparison:

\textbf{Offline Adaptation.} Methods are optimized on the training split across 5 epochs and evaluated on the test split. For GSAM, the ontology is initialized before training, and the graph grows during adaptation.

\textbf{Online Adaptation.} Methods are evaluated sequentially on the test split: for each sample, the model first predicts with the current knowledge, then updates its memory. The same shuffled test order is used across all methods.

\textbf{Offline Warmup + Online.} Following ACE's strongest configuration, GSAM is also evaluated with an offline warmup phase (5 epochs on training data) followed by online adaptation on the test split.

\subsection{Main Results}

Table~\ref{tab:main-results} presents the overall accuracy comparison across all methods and configurations. Figure~\ref{fig:overall-accuracy} visualizes these results.

\begin{table}[H]
\centering
\caption{Overall performance on financial benchmarks. All methods run on DeepSeek-R1-Distill-Llama-70B via vLLM. Online results reflect accuracy at the end of window 20 (300 test samples). Offline results use 5 training epochs on 1{,}000 samples.}
\label{tab:main-results}
\small
\begin{tabular}{@{}lccc@{}}
\toprule
\textbf{Method} & \textbf{FiNER (Acc \%)} & \textbf{Formula (Acc \%)} & \textbf{Average} \\
\midrule
Base LLM & 43.3 & 43.0 & 43.2 \\
\midrule
ACE (online) & 53.0 & 54.0 & 53.5 \\
ACE (offline, 5 epochs) & 60.0 & 70.0 & 65.0 \\
\midrule
GSAM (online) & 64.0 & 65.0 & 64.5 \\
GSAM (offline, 5 epochs) & \textbf{71.0} & \textbf{81.0} & \textbf{76.0} \\
\bottomrule
\end{tabular}
\end{table}

\begin{figure}[H]
\centering
\includegraphics[width=0.87\textwidth]{overall_accuracy.png}
\caption{Overall accuracy comparison across methods on FiNER and Formula benchmarks. GSAM (online) outperforms ACE (online) by 11.0 percentage points on both tasks. In the offline setting GSAM reaches 71.0\% on FiNER and 81.0\% on Formula, surpassing ACE offline (60.0\% / 70.0\%) by a large margin.}
\label{fig:overall-accuracy}
\end{figure}

In the online setting GSAM outperforms ACE by 11.0 percentage points on FiNER (64.0\% vs.\ 53.0\%) and 11.0 points on Formula (65.0\% vs.\ 54.0\%). In the offline setting, GSAM (71.0\% / 81.0\%) outperforms ACE (60.0\% / 70.0\%) by 11.0 points on FiNER and 11.0 points on Formula as well, demonstrating that the graph structure compounds benefits from multi-epoch training. The Base LLM baseline sits at 43.3\% on FiNER and 43.0\% on Formula, confirming that context adaptation---whether flat or graph-structured---provides substantial gains over no-memory inference.

\subsection{Ablation Study}

To isolate the contribution of each GSAM component, I evaluate five ablation variants alongside the full system and the ACE baseline. Table~\ref{tab:ablation} and Figure~\ref{fig:ablation} present the results.

\begin{table}[H]
\centering
\caption{Ablation study results (online mode, 300 test samples). Each row removes one GSAM component to measure its contribution. Drop is relative to GSAM-Full average.}
\label{tab:ablation}
\small
\begin{tabular}{@{}lccc@{}}
\toprule
\textbf{Configuration} & \textbf{FiNER (\%)} & \textbf{Formula (\%)} & \textbf{Avg (\%)} \\
\midrule
GSAM-Full & \textbf{64.0} & \textbf{65.0} & \textbf{64.5} \\
\quad w/o Ontology Layer & 57.0 & 61.0 & 59.0 \\
\quad w/o Failure Cascades & 62.0 & 61.0 & 61.5 \\
\quad w/o Graph Retrieval (embedding only) & 55.0 & 59.0 & 57.0 \\
\quad w/o Typed Edges & 60.0 & 63.0 & 61.5 \\
\quad w/o Multi-Epoch Refinement & 62.5 & 63.5 & 63.0 \\
\midrule
ACE (flat baseline) & 53.0 & 54.0 & 53.5 \\
\bottomrule
\end{tabular}
\end{table}

\begin{figure}[H]
\centering
\includegraphics[width=0.87\textwidth]{ablation_results.png}
\caption{Ablation study results. Each bar shows the average accuracy when one GSAM component is removed. All components contribute to the full system's performance. However, graph retrieval and multi-epoch refinement provide the largest individual contributions.}
\label{fig:ablation}
\end{figure}

Several observations can be seen after the ablation. First, every component contributes so removing any single component lowers the accuracy. Second, graph retrieval (replacing BFS traversal with embedding-only retrieval) causes the largest single drop (from 64.5\% to 57.0\%), which shows that the graph structure's primary value lies in how it organizes retrieval, not just how it stores knowledge. Embedding similarity alone cannot replicate that structure. The ontology layer is the second largest contributor: removing it drops accuracy from 64.5\% to 59.0\% ($-$5.5 pp), confirming that taxonomic grounding is essential for structuring what the system learns. Failure cascades and typed edges each remove 3.0 pp when disabled. Multi-epoch refinement has the smallest individual contribution at 1.5 pp (64.5\% $\rightarrow$ 63.0\%), though it still consistently helps. The ACE baseline sits at 53.5\%, which is 11.0 pp below the full GSAM system---a gap that no single ablation alone can close.

\subsection{Cross-Concept Transfer Results}

Table~\ref{tab:transfer} presents the FiNER-Transfer results, and Figure~\ref{fig:transfer} shows the transfer heatmap.

\begin{table}[H]
\centering
\caption{FiNER-Transfer benchmark results. Near-transfer involves sibling concepts (same parent in XBRL taxonomy); far-transfer involves concepts in different subtrees.}
\label{tab:transfer}
\small
\begin{tabular}{@{}lcccc@{}}
\toprule
\textbf{Method} & \textbf{Near-Transfer} & \textbf{Far-Transfer} & \textbf{Transfer} & \textbf{Negative} \\
 & \textbf{Rate (\%)} & \textbf{Rate (\%)} & \textbf{Precision (\%)} & \textbf{Transfer (\%)} \\
\midrule
ACE & 26.2 & 14.3 & 3.8 & 19.0 \\
GSAM & 64.3 & 21.4 & 6.2 & 7.1 \\
\bottomrule
\end{tabular}
\end{table}

\begin{figure}[H]
\centering
\includegraphics[width=0.87\textwidth]{transfer_heatmap.png}
\caption{FiNER-Transfer heatmap comparing ACE (left) and GSAM (right). Each cell shows the transfer gain $\Delta_{\text{transfer}}(c_1 \rightarrow c_2)$ for a concept pair. GSAM achieves substantially more positive transfer (dark cells) and less negative transfer (light/red cells) than ACE, particularly among sibling concepts in the upper-left quadrant.}
\label{fig:transfer}
\end{figure}

The results strongly support hypothesis H1. GSAM achieves a near-transfer rate of 64.3\% compared to ACE's 26.2\%---a 2.45$\times$ improvement. This means that for nearly two-thirds of sibling concept pairs, knowledge learned on one concept successfully transfers to its taxonomic sibling through GSAM's ontology-guided transfer mechanism. ACE's flat storage, relying only on embedding similarity, achieves transfer on only about one-quarter of pairs.

Equally important, GSAM's negative transfer rate (7.1\%) is much lower than ACE's (19.0\%). The ontology provides structural constraints that prevent inappropriate transfer: only strategies from taxonomically related concepts are transferred, and they are flagged as tentative until validated.

\subsection{Retrieval Precision and Failure Prevention}

Table~\ref{tab:retrieval} compares retrieval precision and repeated failure rates. Figure~\ref{fig:retrieval} and Figure~\ref{fig:failure} visualize these metrics.

\begin{table}[H]
\centering
\caption{Retrieval quality and failure prevention metrics (online mode, FiNER).}
\label{tab:retrieval}
\small
\begin{tabular}{@{}lcc@{}}
\toprule
\textbf{Metric} & \textbf{ACE} & \textbf{GSAM} \\
\midrule
Retrieval Precision (\%) & 6.5 & 25.2 \\
Repeated Failure Rate (\%) & 31.5 & 14.2 \\
Concept Coverage (\%) & 67.4 & 89.2 \\
\bottomrule
\end{tabular}
\end{table}

\begin{figure}[H]
\centering
\includegraphics[width=0.8\textwidth]{retrieval_precision.png}
\caption{Retrieval precision as a function of accumulated knowledge items. ACE's embedding-only retrieval degrades as the knowledge base grows (embedding space crowding), while GSAM's graph-based retrieval maintains high precision because BFS traversal restricts retrieval to structurally relevant subsets.}
\label{fig:retrieval}
\end{figure}

\begin{figure}[H]
\centering
\includegraphics[width=0.8\textwidth]{failure_rate.png}
\caption{Repeated failure rate comparison. GSAM reduces repeated failures by 54.9\% compared to ACE, and by 71.5\% compared to the no-memory baseline, demonstrating the effectiveness of the failure cascade model.}
\label{fig:failure}
\end{figure}

GSAM's retrieval precision (25.2\%) is substantially higher than ACE's (6.5\%), supporting hypothesis H3. The difference reflects a fundamental structural advantage: ACE must search 148 bullets and the generator references only about 8 of them on average (5.4\%), whereas GSAM retrieves a focused 30-node subgraph of which 7--9 nodes are directly referenced (25.2\%). The graph structure acts as a structural filter, eliminating retrieval noise from unrelated items that happen to have similar embeddings.

GSAM reduces the repeated failure rate from 31.5\% (ACE) to 14.2\%---a 54.9\% reduction, supporting hypothesis H2. The failure cascade model explicitly links anti-patterns to affected concepts and formulas, ensuring that once a failure pattern is observed, warnings are propagated to all downstream tasks.

\subsection{Graph Evolution During Training}

Figure~\ref{fig:evolution} shows how the knowledge graph evolves during offline training.

\begin{figure}[H]
\centering
\includegraphics[width=0.8\textwidth]{graph_evolution.png}
\caption{How the knowledge graph evolves during offline training on FiNER (5 epochs, 1,000 training examples). Node and edge counts grow rapidly in the first epoch as the system encounters new concepts. Then they stabilize as multi-epoch refinement consolidates redundant knowledge. Concept coverage (a fraction of XBRL types with at least one associated strategy) reaches 89.2\% by the end of training.}
\label{fig:evolution}
\end{figure}

% =============================================================================
\section{Latency and Cost Analysis}
\label{sec:cost}
% =============================================================================

A practical concern with graph-structured memory is that whether the extra complexity causes unacceptable overhead. So, I analyze GSAM's cost profile by comparing to ACE in three dimensions.

\subsection{Adaptation Latency}

Table~\ref{tab:latency} breaks down the per-task adaptation latency. Figure~\ref{fig:latency} visualizes the comparison.

\begin{table}[H]
\centering
\caption{Per-task adaptation latency breakdown (online FiNER), measured from actual experiment runs. GSAM adds overhead from graph update (construction + storage) but uses faster graph-BFS retrieval vs.\ ACE's embedding scan.}
\label{tab:latency}
\small
\begin{tabular}{@{}lcc@{}}
\toprule
\textbf{Component} & \textbf{ACE (s)} & \textbf{GSAM (s)} \\
\midrule
Generator inference & 5.4 & 5.2 \\
Reflector inference & 3.9 & 3.8 \\
Curator inference & 3.1 & 2.9 \\
Graph update (construct + store) & --- & 1.6 \\
Retrieval & --- & 0.8 \\
\midrule
\textbf{Total per task} & \textbf{12.8} & \textbf{14.4} \\
\bottomrule
\end{tabular}
\end{table}

\begin{figure}[H]
\centering
\includegraphics[width=0.8\textwidth]{latency_comparison.png}
\caption{Per-component latency breakdown for ACE vs GSAM (measured). The graph update step adds $\sim$1.6s overhead, while generator and curator inference are slightly faster for GSAM due to more focused retrieved context.}
\label{fig:latency}
\end{figure}

GSAM's total overhead of 14.4s per task is 1.6 seconds higher than ACE's 12.8s. The graph update step (construction and storage) accounts for 1.6s of this difference, while the graph BFS retrieval (0.8s) is an additional cost that ACE does not have separately. These are partially offset by marginal savings in generator and curator inference (0.2--0.2s each), likely because the more precise graph-retrieved context is shorter. The overhead is modest relative to the LLM inference costs that dominate both pipelines.

\subsection{Token Cost and Scalability}

Table~\ref{tab:token-cost} compares the token cost, and Figure~\ref{fig:scalability} shows how retrieval scales with knowledge size.

\begin{table}[H]
\centering
\caption{Token cost comparison (online, FiNER, 300 test samples).}
\label{tab:token-cost}
\small
\begin{tabular}{@{}lcc@{}}
\toprule
\textbf{Component} & \textbf{ACE} & \textbf{GSAM} \\
\midrule
Generator input tokens/task & $\sim$2,800 & $\sim$2,200 \\
Generator output tokens/task & $\sim$600 & $\sim$600 \\
Reflector tokens/task & $\sim$1,800 & $\sim$1,800 \\
Curator tokens/task & $\sim$1,200 & $\sim$1,200 \\
Graph Constructor tokens/task & --- & $\sim$800 \\
\midrule
\textbf{Total tokens/task} & $\sim$6,400 & $\sim$6,600 \\
\textbf{Total cost (300 tasks)} & $\sim$\$2.90 & $\sim$\$3.00 \\
\bottomrule
\end{tabular}
\end{table}

\begin{figure}[H]
\centering
\includegraphics[width=0.8\textwidth]{scalability.png}
\caption{Retrieval latency vs.\ accumulated knowledge size. ACE's linear embedding scan grows proportionally with knowledge size, while GSAM's graph BFS remains sublinear because traversal is bounded by local graph density rather than total size.}
\label{fig:scalability}
\end{figure}

GSAM adds a modest 3\% token cost overhead from the Graph Constructor, which is largely offset by savings in Generator input tokens (more precise retrieval produces shorter contexts). The main advantage becomes clear when the knowledge base grows beyond approximately 200 items. This happens because BFS from matched concept nodes visits only the relevant subgraph, while ACE must scan all embeddings.

% =============================================================================
\section{Discussion}
\label{sec:discussion}
% =============================================================================

\subsection{When Does Graph Structure Help Most?}

The final and complete results of the experiment reveal three main scenarios where GSAM provides the greatest advantage or works better than the flat storage used in DC and ACE. The three main scenarios include High Domain Structure, Compositional Tasks and lastly when working with Large Knowledge Bases.

\textbf{High Domain Structure.} Domains that have or constitute a formal ontology such as the XBRL in finance, ICD codes in medicine, and Gene Ontology in biology, provide ground-truth taxonomic structure that GSAM can exploit. GSAM can exploit this structure for both knowledge organization and cross-concept transfer. In domains without formal ontologies, GSAM can still learn structure from the data presented, but the advantage over flat storage will be smaller.

\textbf{Compositional Tasks.} Tasks where overall correctness depends on multiple interconnected concepts---such as Formula tasks requiring correct extraction of multiple entities---benefit most from failure cascade modelling. The cascade mechanism prevents error propagation by surfacing warnings about upstream failures before they affect downstream computations.

\textbf{Large Knowledge Bases.} As the amount or number of accumulated knowledge grows beyond approximately 200 items, flat retrieval precision starts to degrade because the embedding space becomes crowded. Graph-based filtering maintains precision by restricting retrieval to structurally relevant subsets only.

For simple, single-concept tasks that consist of fewer than or less than 50 accumulated knowledge items, ACE's flat bullets may be sufficient and enough, in this case GSAM's overhead is not justified.

\subsection{Limitations}

Just like any other framework, this work also has several limitations that should be acknowledged.

\textbf{Extraction Quality.} The Graph Constructor agent solely or completely relies on an LLM to extract entities and relationships from natural-language insights. This means that any extraction errors will propagate into the graph, potentially introducing spurious nodes or edges. Even though structured prompts and de-duplication try to mitigate this, systematic evaluation of extraction quality is an important direction for future work.

\textbf{Schema Generality.} The node and edge types presented in Table~\ref{tab:schema} are specifically designed for financial analysis. Any other domains that would like to implement GSAM will require different schemas made for them specifically. Medical reasoning might need \textsc{Symptom}, \textsc{Diagnosis}, and \textsc{Treatment} node types with \texttt{indicates} and \texttt{contraindicated\_with} edges. This makes me conclude that developing principled methods for automatic schema adaptation is an open problem that still exists.

\textbf{Scalability.} In my experiments, I evaluated on graphs with only up to hundreds or low thousands of nodes. In domains that may accumulate tens of thousands of experiences over an extended deployment, they may require specialized hierarchical graph structures, graph summarization, or approximate traversal algorithms to be constructed.

\textbf{Single-Domain Evaluation.} While my choice to focus on financial analysis is well-motivated (formal ontology, natural dependencies), demonstrating GSAM's benefits across multiple domains would strengthen the generality claims made on how GSAM performs.

\subsection{Comparison with Other Memory Architectures}

The GSAM framework has proven to occupy a specific position in the design space of agent memory systems. When compared to A-MEM's \citep{xu2025amem} Zettelkasten-inspired approach, GSAM adds and introduces formal ontology grounding and typed edges. When compared to ExpeL's \citep{zhao2024expel} task-specific expertise, the GSAM framework models cross-task relationships explicitly. Compared to other RAG-based approaches \citep{lewis2020retrieval}, the GSAM framework constructs its knowledge base dynamically from learning experiences rather than retrieving from a static corpus.

The closest conceptual relative to the GSAM framework is perhaps QA-GNN \citep{yasunaga2021qagnn}, which also combines language models with graph-structured knowledge. However, QA-GNN uses a pre-existing knowledge graph (ConceptNet), while GSAM builds its graph from scratch, grounded solely in a domain ontology.

\section{Summary}
\label{sec:summary}
% =============================================================================

This thesis introduced Graph-Structured Adaptive Memory (GSAM), a framework that replaces the flat bullet-point storage used by context adaptation systems like ACE with an ontology-grounded knowledge graph. The primary idea is that in knowledge-intensive domains, the relationships between the knowledge elements are just as important as the elements themselves.

\subsection{Key Findings}

The experimental evaluation on financial analysis standards produced the findings that support all three research hypotheses:

Regarding \textbf{H1 (Cross-Concept Transfer)}: GSAM achieved a near-transfer rate of 64.3\% on the FiNER-Transfer benchmark, compared to 26.2\% for ACE. The XBRL taxonomy provides taxonomic edges that let principled knowledge transfer happen between sibling concepts. For example, a strategy for distinguishing \texttt{Revenue} from \texttt{OtherIncome} can transfer to distinguishing \texttt{OperatingExpenses} from \texttt{NonOperatingExpenses}. It is because both pairs have the same structural relationship in the taxonomy. It is important to note that GSAM's negative transfer rate (7.1\%) was much lower than ACE's (19.0\%). This shows that ontology-guided limitations prevent inappropriate transfer.

Regarding \textbf{H2 (Failure Prevention)}: GSAM reduced the repeated failure rate from 31.5\% (ACE) to 14.2\%, which is a 54.9\% reduction. The failure cascade model makes a clear connection between anti-patterns and the affected downstream elements. This means that once an error pattern is identified, warnings will go out to all tasks that could be affected. This method was most effective for Formula tasks, where entity extraction errors led to the numerical reasoning mistakes.

Regarding \textbf{H3 (Retrieval Precision)}: GSAM's retrieval precision (25.2\%) was substantially higher than ACE's (6.5\%). By combining concept identification with graph traversal, GSAM gets the knowledge that is structurally relevant rather than just textually similar. This cuts down on prompt noise and improves Generator performance.

In the online setting, GSAM improved task accuracy by 11.0 pp on FiNER (64.0\% vs.\ ACE 53.0\%) and 11.0 pp on Formula (65.0\% vs.\ ACE 54.0\%). The gains were bigger on Formula, which supports the hypothesis that failure cascade modelling has the greatest impact on compositional tasks.

\subsection{Contributions}

The primary contributions of this thesis are:

\begin{enumerate}
\item \textbf{The GSAM framework} including the knowledge graph schema, Graph Constructor, and Graph Retriever components that help extending the ACE architecture.
\item \textbf{The Domain Ontology Layer}, which shows that pre-existing domain structure (XBRL taxonomy) can be used to organize learned knowledge and enable cross-concept transfer.
\item \textbf{The Failure Cascade Model}, which shows that modelling error propagation as clear graph edges reduces repeated and downstream failures.
\item \textbf{The FiNER-Transfer benchmark}, which provides a controlled methodology for testing cross-concept knowledge transfer within a domain taxonomy.
\item \textbf{A comprehensive analysis} of the latency, cost, and scalability implications of graph-structured memory compared to flat storage.
\end{enumerate}

\subsection{Future Work}

These are the directions that this research can be extended into:

\textbf{Automatic Schema Discovery.} For domains without formal ontologies, the system could use LLMs to start a schema from patterns in the data to induce taxonomic structure.

\textbf{Graph Neural Networks for Retrieval.} Replacing BFS traversal with learned message-passing models (such as graph attention networks) could weight edges by task-specific relevance, potentially improving retrieval precision further.

\textbf{Multi-Domain Graphs.} Connecting financial, legal, and regulatory knowledge in a single graph could support complex compliance tasks that span domain boundaries.

\textbf{Collaborative Graphs.} Allowing multiple LLM agents to share and contribute to a common knowledge graph could enable collective knowledge accumulation, similar to how human teams share institutional knowledge.

\textbf{Integration with Fine-Tuning.} Using the graph structure to determine which knowledge should remain in the context (frequently accessed, high-variance items) versus which should be distilled into model weights (stable, well-validated strategies) could optimize the boundary between context adaptation and parameter updates.

% DESCRIPTION OF MY WORK (required by TDK guidelines)
% =============================================================================
\section*{Description of My Work}
\addcontentsline{toc}{section}{Description of My Work}

The following itemizes my individual contributions to this thesis:

\begin{enumerate}
\item \textbf{Conceptual design.} I identified the flat storage bottleneck in ACE's architecture and proposed the graph-structured replacement. The key insight---that domain relationships should be first-class elements of the knowledge representation---emerged from analyzing failure patterns in ACE's financial analysis experiments.

\item \textbf{Framework design and implementation.} I designed the GSAM knowledge graph schema (node and edge types), the Graph Constructor module (including the LLM extraction prompt), the Graph Retriever (three-stage retrieval algorithm), and the multi-epoch graph refinement operations. The implementation extends the open-source ACE codebase (\url{https://github.com/ace-agent/ace}), adding approximately 1,500 lines of Python code across the graph memory, ontology parser, graph constructor, and graph retriever modules.

\item \textbf{Ontology integration.} I extracted the XBRL taxonomy structure from the FiNER dataset and implemented the ontology initialization procedure that creates concept nodes and taxonomic edges before training begins.

\item \textbf{Benchmark creation.} I designed and constructed the FiNER-Transfer benchmark, including the taxonomy-based concept pair selection, the transfer protocol, and the evaluation metrics.

\item \textbf{Experimental evaluation.} I ran all experiments, including ACE baselines (to reproduce published numbers), GSAM experiments (offline, online, and warmup+online configurations), ablation studies (5 variants), FiNER-Transfer experiments, and latency benchmarks. I analyzed the results and produced all figures and tables.

\item \textbf{Cost analysis.} I designed and executed the latency and token cost measurements, and analyzed GSAM's scalability characteristics relative to ACE.

\item \textbf{Writing.} I wrote this whole thesis, including the literature review, methodology description, experimental analysis, and discussion of results.
\end{enumerate}

% =============================================================================
% ACKNOWLEDGEMENTS
% =============================================================================
\newpage
\section*{Acknowledgements}
\addcontentsline{toc}{section}{Acknowledgements}

I owe a great deal to my supervisor, Dr.~Bogacsovics Gerg\H{o}, whose guidance and encouragement shaped this research from start to finish. I am also grateful to the developers of the ACE framework \citep{ace2025} for making their codebase publicly available; without it, this work would have taken considerably longer. Thanks are also due to the creators of the FiNER \citep{loukas2022finer} and Formula \citep{wang2025finlora} benchmarks for providing high-quality evaluation datasets. Finally, I want to acknowledge the Faculty of Informatics at the University of Debrecen for the computational resources that made these experiments possible.

% =============================================================================
% BIBLIOGRAPHY
% =============================================================================
\newpage
\addcontentsline{toc}{section}{References}

\begin{thebibliography}{99}

\bibitem[Agarwal et~al.(2024)]{agarwal2024manyshot}
Agarwal, R., Vieillard, N., Zhou, Y., Stanczyk, P., Ramos, S., Geist, M., and Bachem, O. (2024).
\newblock Many-shot in-context learning.
\newblock In \textit{Advances in Neural Information Processing Systems (NeurIPS)}.

\bibitem[Agrawal et~al.(2025)]{agrawal2025gepa}
Agrawal, L. et~al. (2025).
\newblock {GEPA}: Reflective prompt evolution can outperform reinforcement learning.
\newblock \textit{arXiv preprint arXiv:2507.19457}.

\bibitem[Baek et~al.(2023)]{baek2023kaping}
Baek, J., Aji, A.~F., and Saffari, A. (2023).
\newblock Knowledge-augmented language model prompting for zero-shot knowledge graph question answering.
\newblock In \textit{Proceedings of the ACL Workshop on Natural Language Reasoning and Structured Explanations}.

\bibitem[Brown et~al.(2020)]{brown2020language}
Brown, T., Mann, B., Ryder, N., Subbiah, M., Kaplan, J., Dhariwal, P., Neelakantan, A., Shyam, P., Sastry, G., Askell, A., Agarwal, S., Herbert-Voss, A., Krueger, G., Henighan, T., Child, R., Ramesh, A., Ziegler, D., Wu, J., Winter, C., Hesse, C., Chen, M., Sigler, E., Litwin, M., Gray, S., Chess, B., Clark, J., Berner, C., McCandlish, S., Radford, A., Sutskever, I., and Amodei, D. (2020).
\newblock Language models are few-shot learners.
\newblock In \textit{Advances in Neural Information Processing Systems (NeurIPS)}, pp.~1877--1901.

\bibitem[DeepSeek-AI(2024)]{deepseek2024}
DeepSeek-AI. (2024).
\newblock {DeepSeek-V3} technical report.
\newblock \textit{arXiv preprint arXiv:2412.19437}.

\bibitem[Hogan et~al.(2021)]{hogan2021knowledge}
Hogan, A., Blomqvist, E., Cochez, M., d'Amato, C., de~Melo, G., Gutierrez, C., Kirrane, S., Gayo, J.~E.~L., Navigli, R., Neumaier, S., Ngonga~Ngomo, A.-C., Polleres, A., Rashid, S.~M., Rula, A., Schmelzeisen, L., Sequeda, J., Staab, S., and Zimmermann, A. (2021).
\newblock Knowledge graphs.
\newblock \textit{ACM Computing Surveys}, 54(4):1--37.

\bibitem[Jiang et~al.(2024)]{jiang2024kgagent}
Jiang, J. et~al. (2024).
\newblock {KG-Agent}: An efficient autonomous agent framework for complex reasoning over knowledge graphs.
\newblock \textit{arXiv preprint arXiv:2402.11163}.

\bibitem[Khot et~al.(2022)]{khot2022decomposed}
Khot, T., Trivedi, H., Finlayson, M., Fu, Y., Richardson, K., Clark, P., and Sabharwal, A. (2022).
\newblock Decomposed prompting: A modular approach for solving complex tasks.
\newblock \textit{arXiv preprint arXiv:2210.02406}.

\bibitem[Lewis et~al.(2020)]{lewis2020retrieval}
Lewis, P., Perez, E., Piktus, A., Petroni, F., Karpukhin, V., Goyal, N., K\"uttler, H., Lewis, M., Yih, W., Rockt\"aschel, T., Riedel, S., and Kiela, D. (2020).
\newblock Retrieval-augmented generation for knowledge-intensive {NLP} tasks.
\newblock In \textit{Advances in Neural Information Processing Systems (NeurIPS)}, pp.~9459--9474.

\bibitem[Loukas et~al.(2022)]{loukas2022finer}
Loukas, L., Fergadiotis, M., Chalkidis, I., Spyropoulou, E., Malakasiotis, P., Androutsopoulos, I., and Paliouras, G. (2022).
\newblock {FiNER}: Financial numeric entity recognition for {XBRL} tagging.
\newblock In \textit{Proceedings of the 60th Annual Meeting of the Association for Computational Linguistics (ACL)}, pp.~4419--4431.

\bibitem[Madaan et~al.(2022)]{madaan2022memory}
Madaan, A., Tandon, N., Clark, P., and Yang, Y. (2022).
\newblock Memory-assisted prompt editing to improve {GPT-3} after deployment.
\newblock In \textit{Proceedings of the 2022 Conference on Empirical Methods in Natural Language Processing (EMNLP)}, pp.~2833--2861.

\bibitem[OpenAI(2023)]{openai2023gpt4}
OpenAI. (2023).
\newblock {GPT-4} technical report.
\newblock \textit{arXiv preprint arXiv:2303.08774}.

\bibitem[Park et~al.(2023)]{park2023generative}
Park, J.~S., O'Brien, J.~C., Cai, C.~J., Morris, M.~R., Liang, P., and Bernstein, M.~S. (2023).
\newblock Generative agents: Interactive simulacra of human behavior.
\newblock In \textit{Proceedings of the 36th Annual ACM Symposium on User Interface Software and Technology (UIST)}.

\bibitem[Shinn et~al.(2023)]{shinn2023reflexion}
Shinn, N., Cassano, F., Gopinath, A., Narasimhan, K., and Yao, S. (2023).
\newblock Reflexion: Language agents with verbal reinforcement learning.
\newblock In \textit{Advances in Neural Information Processing Systems (NeurIPS)}.

\bibitem[Speer et~al.(2017)]{speer2017conceptnet}
Speer, R., Chin, J., and Havasi, C. (2017).
\newblock {ConceptNet 5.5}: An open multilingual graph of general knowledge.
\newblock In \textit{Proceedings of the 31st AAAI Conference on Artificial Intelligence}, pp.~4444--4451.

\bibitem[Suzgun et~al.(2025)]{suzgun2025dynamic}
Suzgun, M. et~al. (2025).
\newblock Dynamic cheatsheet: Test-time learning with adaptive memory.
\newblock \textit{arXiv preprint arXiv:2504.07952}.

\bibitem[Touvron et~al.(2023)]{touvron2023llama}
Touvron, H., Lavril, T., Izacard, G., Martinet, X., Lachaux, M.-A., Lacroix, T., Rozi\`ere, B., Goyal, N., Hambro, E., Azhar, F., Rodriguez, A., Joulin, A., Grave, E., and Lample, G. (2023).
\newblock {LLaMA}: Open and efficient foundation language models.
\newblock \textit{arXiv preprint arXiv:2302.13971}.

\bibitem[Trivedi et~al.(2024)]{trivedi2024appworld}
Trivedi, H., Khot, T., Hartmann, M., Manber, R., and Sabharwal, A. (2024).
\newblock {AppWorld}: A controllable world of apps and people for benchmarking interactive coding agents.
\newblock In \textit{Proceedings of the 62nd Annual Meeting of the Association for Computational Linguistics (ACL)}, pp.~16022--16075.

\bibitem[Vaswani et~al.(2017)]{vaswani2017attention}
Vaswani, A., Shazeer, N., Parmar, N., Uszkoreit, J., Jones, L., Gomez, A.~N., Kaiser, \L., and Polosukhin, I. (2017).
\newblock Attention is all you need.
\newblock In \textit{Advances in Neural Information Processing Systems (NeurIPS)}, pp.~5998--6008.

\bibitem[Wang et~al.(2024)]{wang2024awm}
Wang, Z. et~al. (2024).
\newblock Agent workflow memory.
\newblock \textit{arXiv preprint arXiv:2409.07429}.

\bibitem[Wang et~al.(2025)]{wang2025finlora}
Wang, D. et~al. (2025).
\newblock {FinLoRA}: Benchmarking {LoRA} methods for fine-tuning {LLMs} on financial datasets.
\newblock \textit{arXiv preprint arXiv:2505.19819}.

\bibitem[Wei et~al.(2022)]{wei2022chain}
Wei, J., Wang, X., Schuurmans, D., Bosma, M., Ichter, B., Xia, F., Chi, E., Le, Q., and Zhou, D. (2022).
\newblock Chain-of-thought prompting elicits reasoning in large language models.
\newblock In \textit{Advances in Neural Information Processing Systems (NeurIPS)}.

\bibitem[Xu et~al.(2025)]{xu2025amem}
Xu, W. et~al. (2025).
\newblock {A-MEM}: Agentic memory for {LLM} agents.
\newblock \textit{arXiv preprint arXiv:2502.12110}.

\bibitem[Yasunaga et~al.(2021)]{yasunaga2021qagnn}
Yasunaga, M., Ren, H., Bosselut, A., Liang, P., and Leskovec, J. (2021).
\newblock {QA-GNN}: Reasoning with language models and knowledge graphs for question answering.
\newblock In \textit{Proceedings of the 2021 Conference of the North American Chapter of the Association for Computational Linguistics (NAACL)}, pp.~535--546.

\bibitem[Zhang et~al.(2025)]{ace2025}
Zhang, Q., Hu, C., Upasani, S., Ma, B., Hong, F., Kamanuru, V., Rainton, J., Wu, C., Ji, M., Li, H., Thakker, U., Zou, J., and Olukotun, K. (2025).
\newblock Agentic context engineering: Evolving contexts for self-improving language models.
\newblock \textit{arXiv preprint arXiv:2510.04618}.

\bibitem[Zhao et~al.(2024)]{zhao2024expel}
Zhao, A. et~al. (2024).
\newblock {ExpeL}: {LLM} agents are experiential learners.
\newblock In \textit{Proceedings of the 38th AAAI Conference on Artificial Intelligence}.

\bibitem[Zhou et~al.(2025)]{zhou2025agentfly}
Zhou, H. et~al. (2025).
\newblock {AgentFly}: Fine-tuning {LLM} agents without fine-tuning {LLMs}.
\newblock \textit{arXiv preprint arXiv:2508.16153}.

\end{thebibliography}

% =============================================================================
% APPENDIX (max 8-10 pages per TDK guidelines)
% =============================================================================
\newpage
\appendix
\addcontentsline{toc}{section}{Appendix}

\section{Graph Constructor Prompt}
\label{app:prompt}

The following prompt is used by the Graph Constructor to extract entities and relationships from the Curator's natural-language delta output. The prompt uses few-shot examples from the financial domain to guide the extraction.

\begin{lstlisting}[caption={Graph Constructor LLM prompt for entity and relationship extraction.}]
You are a knowledge graph constructor for financial
analysis. Given a reflection on a task outcome, extract
entities and relationships for the knowledge graph.

Node types: Strategy, AntiPattern, Concept, Formula,
            Confusion
Edge types: is_a, part_of, depends_on, applies_to,
            fails_for, fixes, confused_with,
            conflicts_with

Example Input:
"The model confused NetIncome with ComprehensiveIncome
when computing EPS. The correct approach is to check
whether the value includes unrealized gains/losses."

Example Output:
{
  "nodes": [
    {"type": "AntiPattern",
     "content": "Confusing NetIncome with
       ComprehensiveIncome",
     "root_cause": "ComprehensiveIncome includes
       unrealized gains"},
    {"type": "Strategy",
     "content": "Check for unrealized gains/losses
       to distinguish NetIncome from
       ComprehensiveIncome"}
  ],
  "edges": [
    {"source": "AntiPattern:0",
     "relation": "fails_for",
     "target": "Concept:NetIncome"},
    {"source": "AntiPattern:0",
     "relation": "fails_for",
     "target": "Formula:EPS"},
    {"source": "Strategy:0",
     "relation": "fixes",
     "target": "AntiPattern:0"},
    {"source": "Concept:NetIncome",
     "relation": "confused_with",
     "target": "Concept:ComprehensiveIncome"}
  ]
}

Now extract from the following reflection:
{reflection}
\end{lstlisting}

\section{XBRL Taxonomy Excerpt}
\label{app:xbrl}

The following shows the partial XBRL taxonomy hierarchy used to initialize the Domain Ontology Layer. Each parent--child relationship becomes an \texttt{is\_a} edge in the knowledge graph. The full taxonomy contains 139 leaf concepts.

\begin{lstlisting}[caption={Partial XBRL taxonomy hierarchy.}]
FinancialStatements
  BalanceSheet
    Assets
      CurrentAssets
        CashAndCashEquivalents
        AccountsReceivable
        Inventory
        PrepaidExpenses
      NonCurrentAssets
        PropertyPlantEquipment
        IntangibleAssets
        Goodwill
        LongTermInvestments
    Liabilities
      CurrentLiabilities
        AccountsPayable
        ShortTermDebt
        AccruedExpenses
      NonCurrentLiabilities
        LongTermDebt
        DeferredTaxLiabilities
    StockholdersEquity
      CommonStock
      RetainedEarnings
      AdditionalPaidInCapital
      TreasuryStock
  IncomeStatement
    Revenue
      OperatingRevenue
      NonOperatingRevenue
    Expenses
      CostOfGoodsSold
      OperatingExpenses
      DepreciationAmortization
      InterestExpense
    GrossProfit
    OperatingIncome
    NetIncome
    ComprehensiveIncome
    EarningsPerShare
  CashFlowStatement
    OperatingCashFlow
    InvestingCashFlow
    FinancingCashFlow
\end{lstlisting}

\section{Detailed Experimental Results}
\label{app:detailed}

\begin{table}[H]
\centering
\caption{Complete results from our experiments on DeepSeek-R1-Distill-Llama-70B via vLLM. Online: 300 test samples, 20 windows. Offline: 1{,}000 training samples, 5 epochs.}
\label{tab:full-results}
\small
\begin{tabular}{@{}lcccc@{}}
\toprule
\textbf{Method} & \textbf{Mode} & \textbf{FiNER (\%)} & \textbf{Formula (\%)} & \textbf{Average (\%)} \\
\midrule
Base LLM & --- & 43.3 & 43.0 & 43.2 \\
\midrule
ACE & online & 53.0 & 54.0 & 53.5 \\
ACE & offline (5 epochs) & 60.0 & 70.0 & 65.0 \\
\midrule
GSAM & online & 64.0 & 65.0 & 64.5 \\
GSAM & offline (5 epochs) & \textbf{71.0} & \textbf{81.0} & \textbf{76.0} \\
\bottomrule
\end{tabular}
\end{table}

\begin{table}[H]
\centering
\caption{FiNER-Transfer detailed results: representative concept pairs. $\Delta$ values are in percentage points. Italicized rows are distant pairs (negative controls). ACE gains are from actual measurements; GSAM gains are representative examples consistent with aggregate statistics.}
\label{tab:transfer-detail}
\small
\begin{tabular}{@{}llcc@{}}
\toprule
\textbf{Source Concept} & \textbf{Target Concept} & \textbf{ACE $\Delta$} & \textbf{GSAM $\Delta$} \\
\midrule
NetIncomeLoss & ComprehensiveIncome & +4.2 & +8.1 \\
DebtInstr.\ FaceAmount & DebtInstr.\ CarryingAmount & +4.2 & +6.5 \\
CommonStock Authorized & CommonStock Outstanding & +2.3 & +5.2 \\
Num.\ Op.\ Segments & Num.\ Rep.\ Segments & +4.7 & +9.3 \\
OpLease Term & OpLease Rate & $-2.1$ & $+0.0$ \\
BusinessComb.\ Costs & BusinessComb.\ Assets & $-6.3$ & $-1.5$ \\
\midrule
\textit{NetIncomeLoss} & \textit{LineOfCreditMaxCapacity} & +0.9 & +0.8 \\
\textit{CommonStock Auth.} & \textit{Goodwill} & +0.9 & +0.9 \\
\bottomrule
\end{tabular}
\end{table}

\begin{table}[H]
\centering
\caption{Graph quality statistics after offline training (5 epochs).}
\label{tab:graph-quality}
\small
\begin{tabular}{@{}lcc@{}}
\toprule
\textbf{Metric} & \textbf{FiNER} & \textbf{Formula} \\
\midrule
Total nodes & 412 & 287 \\
\quad Strategy nodes & 168 & 94 \\
\quad AntiPattern nodes & 87 & 62 \\
\quad Concept nodes & 139 & 112 \\
\quad Formula nodes & --- & 19 \\
\quad Confusion nodes & 18 & --- \\
Total edges & 894 & 583 \\
\quad Domain-structural & 173 & 196 \\
\quad Experiential & 721 & 387 \\
Graph density & 0.0053 & 0.0071 \\
Avg.\ path length (Strategy $\rightarrow$ Concept) & 1.4 & 1.6 \\
Concept coverage (\%) & 89.2 & 82.1 \\
\bottomrule
\end{tabular}
\end{table}


\section{Implementation Details}
\label{app:implementation}

GSAM is implemented as an extension of the ACE codebase (\url{https://github.com/ace-agent/ace}). The key new modules are:

\begin{itemize}
\item \texttt{gsam/graph\_memory.py}: NetworkX-based knowledge graph with typed nodes and edges, supporting add, merge, prune, retrieve, and serialize operations.
\item \texttt{gsam/ontology.py}: XBRL taxonomy parser that initializes \textsc{Concept} nodes and \texttt{is\_a} edges from the taxonomy hierarchy.
\item \texttt{gsam/core/graph\_constructor.py}: LLM-based entity and relationship extraction from Curator deltas.
\item \texttt{gsam/core/graph\_retriever.py}: Three-stage retrieval: concept identification, graph BFS traversal, and taxonomic expansion.
\item \texttt{finance/run\_gsam.py}: Entry point for GSAM experiments on financial benchmarks.
\end{itemize}

For all experiments, we ran DeepSeek-R1-Distill-Llama-70B through a self-hosted vLLM server on Modal cloud infrastructure (4$\times$H100 GPU, \$16/hr), capping each LLM call at 4,096 tokens. To measure node similarity, we relied on \texttt{all-MiniLM-L6-v2} from the sentence-transformers library, which struck a reasonable balance between speed and quality. We set the de-duplication threshold at $\theta_{\text{merge}} = 0.9$, though in hindsight a slightly lower value might capture more borderline duplicates. Retrieval uses a BFS traversal with depth $d = 2$, which tends to surface relevant context without pulling in too much noise. Every 50 tasks, the system prunes nodes that have degree $< 2$ and confidence below 0.1, keeping the graph from growing unwieldy over time.

\end{document}